\begin{theorem}[Fixed-resolvent defect identity for the Weyl family]
    \label{thm:weyl_resolvent_defect_identity}
    \lean{Matrix.weyl_resolvent_defect_identity}
    \leanok
    \uses{thm:weyl_operator_unitary}
    Let $\rho$ and $\sigma$ be positive definite on
    $\mathcal H_S\otimes\mathbb C^{d_C}$, let
    $q=d_C^{-2}$, and put
    \[
        A_g=qU_g\rho U_g^{\dagger},\qquad
        B_g=qU_g\sigma U_g^{\dagger},\qquad
        A=\sum_g A_g,\qquad B=\sum_g B_g,
    \]
    where $U_g=\mathbf 1_S\otimes W_g$.  For $t>0$, let
    \[
        T_g(X)=A_gX+tXB_g,\qquad T(X)=AX+tXB,
        \qquad \Lambda_t=T^{-1}(B).
    \]
    Then
    \[
        \sum_g\left\|
          T_g^{-1/2}(B_g)-T_g^{1/2}(\Lambda_t)
        \right\|_{\mathrm{HS}}^2
        =\sum_g\re
          \langle B_g,T_g^{-1}(B_g)\rangle_{\mathrm{HS}}
          -\re
          \langle B,T^{-1}(B)\rangle_{\mathrm{HS}}.
    \]
    This is the positive-definite, fixed-$t$ identity in
    Jen\v{c}ov\'a--Ruskai, arXiv:0903.2895v4, Appendix,
    lines 1313--1343.  It does not infer zero defect from equality of relative
    entropies and makes no assertion about singular supports.
\end{theorem}

\begin{proof}\leanok
    Write $S_g$ for the positive definite matrix representing $T_g$ under
    the vectorization $X\mapsto\opvec(X^{\mathsf T})$, and put
    $b_g=\opvec(B_g^{\mathsf T})$ and
    $x=S^{-1}\sum_g b_g$.  For the residual $r_g=b_g-S_gx$, direct expansion
    gives
    \[
      \sum_g\langle r_g,S_g^{-1}r_g\rangle
        =\sum_g\langle b_g,S_g^{-1}b_g\rangle
          -\left\langle\sum_gb_g,
            S^{-1}\sum_gb_g\right\rangle.
    \]
    Since $S_g^{1/2}$ is invertible,
    \[
      \|S_g^{-1/2}b_g-S_g^{1/2}x\|^2
        =\|S_g^{-1/2}r_g\|^2
        =\re\langle r_g,S_g^{-1}r_g\rangle.
    \]
    Summing proves the identity.
\end{proof}

\begin{theorem}[Source-$A$ fixed-resolvent defect identity]
    \label{thm:weyl_sourceA_resolvent_defect_identity}
    \lean{Matrix.weyl_sourceA_resolvent_defect_identity}
    \leanok
    \uses{thm:weyl_operator_unitary}
    Under the notation of Theorem~\ref{thm:weyl_resolvent_defect_identity},
    set $\Gamma_t=T^{-1}(A)$.  Then, for every $t>0$,
    \[
      \sum_g\left\|
        T_g^{-1/2}(A_g)-T_g^{1/2}(\Gamma_t)
      \right\|_{\rm HS}^2
      =\sum_g\re
        \langle A_g,T_g^{-1}(A_g)\rangle_{\rm HS}
       -\re\langle A,T^{-1}(A)\rangle_{\rm HS}.
    \]
    This is the second defect family in Jen\v{c}ov\'a--Ruskai,
    arXiv:0903.2895v4, \S4 and Appendix.
\end{theorem}

\begin{proof}\leanok
    Apply the residual identity of
    Theorem~\ref{thm:weyl_resolvent_defect_identity} with
    $a_g=\opvec(A_g^{\mathsf T})$ and
    $\bar a=\opvec(A^{\mathsf T})$.  Thus
    \[
      \sum_g\langle a_g,S_g^{-1}a_g\rangle
        -\langle\bar a,S^{-1}\bar a\rangle
      =\sum_g\|S_g^{-1/2}a_g-S_g^{1/2}S^{-1}\bar a\|^2,
    \]
    which is the asserted source-$A$ identity.
\end{proof}

\begin{theorem}[Finite-Weyl relative-entropy gap integral]
    \label{thm:weyl_relative_entropy_gap_integral}
    \lean{Matrix.weightedWeylConjugate}
    \lean{Matrix.weightedWeylAverage}
    \lean{Matrix.weylSourceAResolventDefect}
    \lean{Matrix.weylSourceBResolventDefect}
    \lean{Matrix.weylRelativeEntropyGap}
    \lean{Matrix.weyl_relativeEntropy_gap_integral}
    \lean{Matrix.weylSourceAResolventDefect_continuousOn}
    \lean{Matrix.weylSourceBResolventDefect_continuousOn}
    \lean{Matrix.weylSourceAResolventDefect_nonneg}
    \lean{Matrix.weylSourceBResolventDefect_nonneg}
    \lean{Matrix.weylRelativeEntropyGapIntegrand_integrableOn}
    \lean{Matrix.weylRelativeEntropyGapIntegrand_continuousOn}
    \lean{Matrix.weylRelativeEntropyGapIntegrand_nonneg}
    \leanok
    \uses{thm:relative_entropy_resolvent_integral,
        thm:weyl_resolvent_defect_identity,
        thm:weyl_sourceA_resolvent_defect_identity}
    For the positive definite finite-Weyl family, let
    \[
      \Def_A(t)=
        \sum_g\re\langle A_g,T_g^{-1}(A_g)\rangle_{\rm HS}
        -\re\langle A,T^{-1}(A)\rangle_{\rm HS}
    \]
    and define $\Def_B(t)$ analogously, with the real parts of
    the corresponding source-$B$ pairings.  Then both defects are
    nonnegative and continuous for $t>0$, and
    \[
      \sum_gD(A_g\Vert B_g)-D(A\Vert B)
      =\int_0^\infty
        \frac{\Def_A(t)+t\Def_B(t)}{1+t}\,dt.
    \]
    The integrand is continuous, nonnegative, and integrable on
    $(0,\infty)$.  The coefficient of the source-$B$ defect is exactly $t$,
    as prescribed by the integral formula and equality analysis of
    Jen\v{c}ov\'a--Ruskai, arXiv:0903.2895v4, \S4.
\end{theorem}

\begin{proof}\leanok
    Apply Theorem~\ref{thm:relative_entropy_resolvent_integral} to each pair
    $(A_g,B_g)$ and to $(A,B)$, and interchange the finite sum with the
    integral.  The trace terms cancel because $B=\sum_gB_g$.  The two
    fixed-resolvent identities write the defects as sums of squared norms,
    proving nonnegativity.  Continuity and integrability follow from the
    corresponding spectral assertions before taking the finite difference.
\end{proof}

\begin{theorem}[Zero Weyl defect gives the common resolvent solution]
    \label{thm:weyl_resolvent_zero_defect}
    \lean{Matrix.weyl_resolvent_eq_of_defect_eq_zero}
    \leanok
    \uses{thm:weyl_resolvent_defect_identity}
    Under the hypotheses and notation of the preceding theorem, suppose that
    \begin{equation}\label{eq:weyl_resolvent_zero_defect_hyp}
        \sum_g\re
          \langle B_g,T_g^{-1}(B_g)\rangle_{\mathrm{HS}}
          -\re
          \langle B,T^{-1}(B)\rangle_{\mathrm{HS}}=0.
    \end{equation}
    Then, for every Weyl index $g$,
    \[
        T_g^{-1}(B_g)=T^{-1}(B).
    \]
    In particular this holds for the identity Weyl element $g=(0,0)$.
    This is the common-resolvent conclusion in Jen\v{c}ov\'a--Ruskai,
    arXiv:0903.2895v4, \S4, lines 652--674; its squared-defect input is in
    Appendix, lines 1313--1343.
    This fixed-$t$, positive-definite conclusion does not assert that equality
    of relative entropies implies the displayed scalar hypothesis.
\end{theorem}

\begin{proof}\leanok
    Equation~(\ref{eq:weyl_resolvent_zero_defect_hyp}) writes the hypothesis
    as a finite sum of squared norms.  Each term is nonnegative, so every term
    vanishes.  Thus
    $S_g^{-1/2}(b_g-S_gx)=0$ for every $g$.  Invertibility of
    $S_g^{-1/2}$ gives $b_g=S_gx$, and hence
    $S_g^{-1}b_g=x=S^{-1}\sum_hb_h$.
\end{proof}

\begin{theorem}[Zero finite-Weyl gap gives the common resolvent solution]
    \label{thm:weyl_relative_entropy_zero_gap_resolvent}
    \lean{Matrix.weyl_sourceB_defect_eq_zero_of_gap_eq_zero}
    \lean{Matrix.weyl_resolvent_eq_of_relativeEntropy_gap_eq_zero}
    \lean{Matrix.weyl_resolvent_eq_of_partialTraceRight_eq}
    \leanok
    \uses{thm:partial_trace_saturation_weyl_gap_zero,
        thm:weyl_relative_entropy_gap_integral,
        thm:weyl_resolvent_zero_defect}
    Under the positive-definite finite-Weyl hypotheses, suppose that
    \[
      \sum_gD(A_g\Vert B_g)-D(A\Vert B)=0.
    \]
    Then $\Def_B(t)=0$ for every $t>0$, and consequently
    \[
      T_g^{-1}(B_g)=T^{-1}(B)
    \]
    for every $t>0$ and every Weyl index $g$.
    In particular, equality of relative entropy under the right partial trace
    implies this conclusion by
    Theorem~\ref{thm:partial_trace_saturation_weyl_gap_zero}.
    This is the positive-definite conclusion of the equality argument in
    Jen\v{c}ov\'a--Ruskai, arXiv:0903.2895v4, \S4 and Appendix.  It makes no
    assertion at $t=0$ or for singular inputs.
\end{theorem}

\begin{proof}\leanok
    The integrand in
    Theorem~\ref{thm:weyl_relative_entropy_gap_integral} is nonnegative and
    has integral zero, hence it vanishes almost everywhere.  Its continuity
    improves this to vanishing at every $t>0$.  Since
    $\Def_A(t)\geq0$, $\Def_B(t)\geq0$, and
    $t/(1+t)>0$, it follows that $\Def_B(t)=0$.  Theorem
    \ref{thm:weyl_resolvent_zero_defect} now gives the common solution.
\end{proof}

\begin{lemma}[Positive square root from common resolvents]
    \label{lem:positive_sqrt_from_common_resolvents}
    \lean{Matrix.sqrt_mulVec_eq_of_resolvent_mulVec_eq}
    \lean{Matrix.mulVec_sqrt_mulVec_eq_of_mulVec_resolvent_mulVec_eq}
    \lean{Matrix.sourceB_resolvent_eq_relativeModular}
    \lean{Matrix.relativeModular_sqrt_mulVec_vec_one}
    \leanok
    Let $S,T$ be positive semidefinite matrices and let $x$ be a vector.  If,
    for every $t>0$,
    \[
        (t\mathbf 1+S)^{-1}x=(t\mathbf 1+T)^{-1}x,
    \]
    then $\sqrt S\,x=\sqrt T\,x$.
    More generally, for any fixed matrix $Q$, if
    \[
        Q(t\mathbf 1+S)^{-1}x=Q(t\mathbf 1+T)^{-1}x
    \]
    for every $t>0$, then $Q\sqrt S\,x=Q\sqrt T\,x$.

    In particular, for positive definite $A,B$ and every $t>0$, put
    $\Delta_{A,B}=A\otimes(B^{-1})^{\mathsf T}$.  The source-$B$ left--right
    resolvent satisfies
    \[
      \left(A\otimes\mathbf 1+t\,\mathbf 1\otimes B^{\mathsf T}\right)^{-1}
        \opvec(B^{\mathsf T})
      =\left(t\mathbf 1+\Delta_{A,B}\right)^{-1}
        \opvec(\mathbf 1^{\mathsf T}),
    \]
    and
    \[
      \sqrt{\Delta_{A,B}}\,\opvec(\mathbf 1^{\mathsf T})
        =\opvec\!\left(
          (\sqrt A\,(\sqrt B)^{-1})^{\mathsf T}\right).
    \]
    These are the positive-square-root specializations of the passage from
    relative modular resolvents to analytic functions of the relative modular
    operator in Jen\v{c}ov\'a--Ruskai, arXiv:0903.2895v4, lines 658--680.
\end{lemma}

\begin{proof}\leanok
    Use the L\"owner integral representation of the power $p=1/2$.  Its
    integrand at $t>0$ is
    \[
      f_t(S)=t^{-1/2}\mathbf 1-t^{1/2}(t\mathbf 1+S)^{-1}.
    \]
    Applied to $x$, this is
    \[
      f_t(S)x=t^{-1/2}x-t^{1/2}(t\mathbf 1+S)^{-1}x.
    \]
    The hypothesis therefore gives
    $Q\bigl(f_t(S)x\bigr)=Q\bigl(f_t(T)x\bigr)$ for every $t>0$.
    Since $M\mapsto Q(Mx)$ is a bounded linear map from
    $\mathbb{M}_n(\mathbb{C})$ to $\mathbb{C}^n$,
    \[
      Q\left(\left(\int_0^\infty f_t(S)\,d\mu(t)\right)x\right)
        =\int_0^\infty Q\bigl(f_t(S)x\bigr)\,d\mu(t),
    \]
    and similarly for $T$.  Integration therefore gives
    $Q\sqrt S\,x=Q\sqrt T\,x$.

    For the second assertion, factor
    \[
      L_A+tR_B=(\Delta_{A,B}+t\mathbf 1)R_B.
    \]
    Applying the inverse to $B=R_B(\mathbf 1)$ gives the shifted relative
    modular resolvent on $\mathbf 1$.  Finally,
    \[
      \sqrt{\Delta_{A,B}}
        =\sqrt A\otimes\bigl((\sqrt B)^{-1}\bigr)^{\mathsf T},
    \]
    as follows from uniqueness of the positive square root.
\end{proof}

\begin{lemma}[Right multiplication under transposed vectorization]
    \label{lem:right_multiplication_transposed_vectorization}
    \lean{Matrix.one_kronecker_transpose_mulVec_vec_transpose}
    \leanok
    For matrices $M$ and $P$, column-stacking vectorization satisfies
    \[
      \left(\mathbf 1\otimes P^{\mathsf T}\right)
        \opvec(M^{\mathsf T})
      =\opvec\!\left((MP)^{\mathsf T}\right).
    \]
\end{lemma}

\begin{proof}\leanok
    This is the Kronecker vectorization identity
    $(B\otimes A)\opvec(X)=\opvec(AXB^{\mathsf T})$
    with $A=P^{\mathsf T}$, $B=\mathbf 1$, and $X=M^{\mathsf T}$.
\end{proof}

\begin{lemma}[Canonical support-domain source resolvent solution]
    \label{lem:support_relative_modular_source_solution}
    \lean{Matrix.supportRelativeModular_sourceB_solution}
    \leanok
    \uses{def:support_inverse_square_root}
    Let $A,B$ be positive semidefinite, let $t>0$, set
    $B^+=(B^{-1/2}_{\mathrm{supp}})^2$, and let $P_B$ be the support
    projection of $B$.  Then
    \[
      \left(A\otimes\mathbf 1+t(\mathbf 1\otimes B^{\mathsf T})\right)
      (\mathbf 1\otimes P_B^{\mathsf T})
      \left(t\mathbf 1+A\otimes(B^+)^{\mathsf T}\right)^{-1}
      \opvec(\mathbf 1^{\mathsf T})
      =\opvec(B^{\mathsf T}).
    \]
\end{lemma}

\begin{proof}\leanok
    \uses{lem:support_inverse_square_root_generalized_inverse,
      lem:right_multiplication_transposed_vectorization}
    Put $C=\mathbf 1\otimes B^{\mathsf T}$.  The support generalized-inverse
    identity gives
    \[
      C\left(A\otimes(B^+)^{\mathsf T}\right)
        =A\otimes P_B^{\mathsf T}.
    \]
    Together with $B^{\mathsf T}P_B^{\mathsf T}=B^{\mathsf T}$, this factors
    the left two operators as
    \[
      \left(A\otimes\mathbf 1+tC\right)
        (\mathbf 1\otimes P_B^{\mathsf T})
      =C\left(t\mathbf 1+A\otimes(B^+)^{\mathsf T}\right).
    \]
    The shifted relative-modular matrix is positive definite and hence
    invertible, so
    \[
      \left(A\otimes\mathbf 1+tC\right)
        (\mathbf 1\otimes P_B^{\mathsf T})
        \left(t\mathbf 1+A\otimes(B^+)^{\mathsf T}\right)^{-1}
      =C.
    \]
    The Kronecker vectorization identity gives
    \[
      C\opvec(\mathbf 1^{\mathsf T})
        =\opvec(B^{\mathsf T}),
    \]
    which is the claim.
\end{proof}

\begin{lemma}[Support relative-modular square-root evaluation]
    \label{lem:support_relative_modular_sqrt_evaluation}
    \lean{Matrix.supportRelativeModular_sqrt_mulVec_vec_one}
    \leanok
    \uses{def:support_inverse_square_root}
    Let $A$ and $B$ be positive semidefinite, and write
    $B^+=(B^{-1/2}_{\mathrm{supp}})^2$.  Then
    \[
      \sqrt{A\otimes(B^+)^{\mathsf T}}\,
        \opvec(\mathbf 1^{\mathsf T})
      =\opvec\!\left(
        (\sqrt A\,B^{-1/2}_{\mathrm{supp}})^{\mathsf T}\right).
    \]
\end{lemma}

\begin{proof}\leanok
    \uses{lem:support_inverse_square_root_positive_semidefinite}
    The support inverse square root is positive semidefinite, and
    \[
      \left(\sqrt A\otimes
        (B^{-1/2}_{\mathrm{supp}})^{\mathsf T}\right)^2
      =A\otimes(B^+)^{\mathsf T}.
    \]
    Uniqueness of the positive square root gives the corresponding Kronecker
    factorization.  Applying the Kronecker vectorization identity to
    $\opvec(\mathbf 1^{\mathsf T})$ gives the stated equation.
\end{proof}

\begin{lemma}[Relative modular square-root ratio on the support]
    \label{lem:support_relative_modular_sqrt_ratio}
    \lean{Matrix.supportRelativeModular_sqrt_ratio_eq_of_resolvent_mulVec_eq}
    \leanok
    \uses{def:support_inverse_square_root}
    Let $A,B,C,D$ be positive semidefinite matrices.  Write
    $B^+=(B^{-1/2}_{\mathrm{supp}})^2$ and
    $D^+=(D^{-1/2}_{\mathrm{supp}})^2$, and let $P_B$ be the orthogonal
    projection onto $(\ker B)^\perp$.  If, for every $t>0$,
    \[
      \left[
        \left(t\mathbf 1+A\otimes(B^+)^{\mathsf T}\right)^{-1}
          (\mathbf 1)
      \right]P_B
      =
      \left[
        \left(t\mathbf 1+C\otimes(D^+)^{\mathsf T}\right)^{-1}
          (\mathbf 1)
      \right]P_B,
    \]
    where the inverses act as relative-modular superoperators on matrices, then
    \[
      \left(\sqrt A\,B^{-1/2}_{\mathrm{supp}}\right)P_B
        =\left(\sqrt C\,D^{-1/2}_{\mathrm{supp}}\right)P_B.
    \]
    This is the square-root specialization of the support functional-calculus
    passage in Jen\v{c}ov\'a--Ruskai, arXiv:0903.2895v4, lines 788--793.  The
    projection $P_B$ is essential: the source gives the common generalized
    resolvents only after restriction to $(\ker B)^\perp$.  Deriving this
    restricted equality requires the preceding singular equality argument and
    its kernel hypotheses.
    % Formalization status: the preceding equality-to-resolvent step remains open.
\end{lemma}

\begin{proof}\leanok
    \uses{lem:positive_sqrt_from_common_resolvents,
        lem:support_relative_modular_sqrt_evaluation,
        lem:right_multiplication_transposed_vectorization}
    By Lemma~\ref{lem:support_relative_modular_sqrt_evaluation},
    \[
      \sqrt{A\otimes(B^+)^{\mathsf T}}\,
        \opvec(\mathbf 1^{\mathsf T})
      =\opvec\!\left(
        (\sqrt A\,B^{-1/2}_{\mathrm{supp}})^{\mathsf T}\right).
    \]
    Setting $Q=\mathbf 1\otimes P_B^{\mathsf T}$ and applying the common-
    resolvent lemma (Lemma~\ref{lem:positive_sqrt_from_common_resolvents}) gives
    \[
      \opvec\!\left(
        (\sqrt A\,B^{-1/2}_{\mathrm{supp}}P_B)^{\mathsf T}\right)
      =\opvec\!\left(
        (\sqrt C\,D^{-1/2}_{\mathrm{supp}}P_B)^{\mathsf T}\right).
    \]
    Since
    \[
      \opvec(X^{\mathsf T})
        =\opvec(Y^{\mathsf T})\quad\Longrightarrow\quad X=Y,
    \]
    injectivity of vectorization gives the conclusion.
\end{proof}

\begin{theorem}[Positive-definite saturation gives the identity Weyl sandwich]
    \label{thm:partial_trace_saturation_identity_weyl_sandwich}
    \lean{Matrix.weyl_sqrt_ratio_eq_of_partialTraceRight_eq}
    \lean{Matrix.weyl_identity_sandwich_of_partialTraceRight_eq}
    \lean{Matrix.partialTraceRightPetzMap_eq_of_relativeEntropy_eq_posDef}
    \leanok
    Let $\rho$ and $\sigma$ be positive definite and suppose that
    \[
      D(\rho\Vert\sigma)
        =D(\tr_C\rho\Vert\tr_C\sigma).
    \]
    Put
    \[
      \overline\rho=(\tr_C\rho)\otimes d_C^{-1}\mathbf 1_C,
      \qquad
      \overline\sigma=(\tr_C\sigma)\otimes d_C^{-1}\mathbf 1_C.
    \]
    Then
    \[
      \sqrt\rho\,\sigma^{-1/2}
        =\sqrt{\overline\rho}\,\overline\sigma^{-1/2},
    \]
    and consequently
    \[
      \sqrt\sigma\,\overline\sigma^{-1/2}\,
      \overline\rho\,\overline\sigma^{-1/2}\,\sqrt\sigma=\rho.
    \]
    Hence the raw partial-trace Petz map satisfies
    \[
      \mathcal R_\sigma(\tr_C\rho)=\rho.
    \]
    This is the positive-definite case only; no singular-support conclusion is
    asserted.
\end{theorem}

\begin{proof}\leanok
    \uses{thm:weyl_relative_entropy_zero_gap_resolvent,
        lem:positive_sqrt_from_common_resolvents,
        thm:weyl_identity_sandwich_petz_recovery}
    Apply the common-resolvent theorem to the identity Weyl summand and the
    Weyl average.  The preceding lemma converts the result to equality of the
    two square-root ratios.  For $c=d_C^{-2}$, the uniform scalar in the
    identity summand cancels through
    \[
      \sqrt{c\rho}\,(\sqrt{c\sigma})^{-1}
      =(\sqrt c\,\sqrt\rho)
        \bigl((\sqrt c)^{-1}(\sqrt\sigma)^{-1}\bigr)
      =\sqrt\rho\,(\sqrt\sigma)^{-1}.
    \]
    The Weyl twirl identifies the average with the displayed maximally mixed
    extensions.

    Taking the adjoint product of the ratio equality yields
    \[
      \sigma^{-1/2}\rho\sigma^{-1/2}
        =\overline\sigma^{-1/2}\overline\rho
          \overline\sigma^{-1/2}.
    \]
    Multiplication by $\sqrt\sigma$ on both sides gives the sandwich.  The raw
    Petz recovery identity follows from
    Theorem~\ref{thm:weyl_identity_sandwich_petz_recovery}.
\end{proof}

\begin{lemma}[Maximally mixed support-sandwich normalization]
    \label{lem:maximally_mixed_support_sandwich}
    \lean{Matrix.PosSemidef.supportInvSqrt_kronecker_maximallyMixed_mul_mul}
    \leanok
    \uses{lem:support_inverse_square_root_maximally_mixed}
    Let $\tau$ be positive semidefinite on $\mathcal H_S$, let $X$ be a
    matrix on $\mathcal H_S$, and let
    $\overline\tau=\tau\otimes d_C^{-1}\mathbf 1_C$.  Then
    \[
        \overline\tau^{-1/2}_{\mathrm{supp}}
        \bigl(X\otimes d_C^{-1}\mathbf 1_C\bigr)
        \overline\tau^{-1/2}_{\mathrm{supp}}
        =\bigl(\tau^{-1/2}_{\mathrm{supp}}X
          \tau^{-1/2}_{\mathrm{supp}}\bigr)\otimes\mathbf 1_C.
    \]
\end{lemma}

\begin{proof}\leanok
    By the preceding support inverse formula, each outer factor contributes
    $\sqrt{d_C}$.  The scalar coefficient cancels because
    \[
        \sqrt{d_C}\,d_C^{-1}\sqrt{d_C}
          =d_C^{-1}(\sqrt{d_C})^2=1.
    \]
    The remaining matrix product is the asserted unital tensor embedding.
\end{proof}

\begin{theorem}[Identity Weyl sandwich implies raw Petz recovery]
    \label{thm:weyl_identity_sandwich_petz_recovery}
    \lean{Matrix.partialTraceRightPetzMap_eq_of_weyl_identity_sandwich}
    \leanok
    \uses{def:partial_trace_support_petz_map,
        lem:partial_trace_support_petz_map_apply,
        lem:maximally_mixed_support_sandwich}
    Let $\rho$ and $\sigma$ be matrices on
    $\mathcal H_S\otimes\mathbb C^{d_C}$, with $\sigma$ positive semidefinite,
    and set
    \[
        \overline\sigma=(\tr_C\sigma)\otimes d_C^{-1}\mathbf 1_C.
    \]
    If the identity summand obeys the support sandwich identity
    \[
        \sqrt\sigma\,
        \overline\sigma^{-1/2}_{\mathrm{supp}}
        \bigl((\tr_C\rho)\otimes d_C^{-1}\mathbf 1_C\bigr)
        \overline\sigma^{-1/2}_{\mathrm{supp}}\,
        \sqrt\sigma=\rho,
    \]
    then the raw support Petz map recovers $\rho$:
    \[
        \mathcal R_\sigma(\tr_C\rho)=\rho.
    \]
    This is the algebraic reduction in
    \cite[Theorem~3, equation~(8)]{Hayden2004Structure}.  It does not derive
    the displayed sandwich identity from equality of relative entropies.
\end{theorem}

\begin{proof}\leanok
    The preceding normalization rewrites the middle three factors as
    \[
        \bigl((\tr_C\sigma)^{-1/2}_{\mathrm{supp}}
          (\tr_C\rho)
          (\tr_C\sigma)^{-1/2}_{\mathrm{supp}}\bigr)
          \otimes\mathbf 1_C.
    \]
    The support Petz formula therefore identifies the left-hand side of the
    assumed sandwich identity with $\mathcal R_\sigma(\tr_C\rho)$.
\end{proof}

\begin{theorem}[Relative entropy against the product of the marginals]
    \label{thm:relative_entropy_product_marginals}
    \lean{quantumRelativeEntropy_product_marginals}
    \leanok
    \uses{def:quantum_relative_entropy, def:von_neumann_entropy}
    For every positive semidefinite operator $\omega_{XY}$,
    \[
        D(\omega_{XY}\,\|\,\omega_X\otimes\omega_Y)
        =S(\omega_X)+S(\omega_Y)-S(\omega_{XY}).
    \]
    No invertibility assumption is made on either marginal.
    This is \cite[Equation~(4)]{Hayden2004Structure}.
\end{theorem}

\begin{proof}\leanok
    \uses{lem:log_kronecker_possemidef,
        thm:marginal_support_left_absorption,
        thm:right_marginal_support_left_absorption,
        lem:trace_lift_left_factor_mul, lem:trace_lift_right_factor_mul}
    The lifted support projections $P_X\otimes\mathbf 1_Y$ and
    $\mathbf 1_X\otimes P_Y$ both fix $\omega_{XY}$.  Substituting
    Lemma~\ref{lem:log_kronecker_possemidef} into the cross term of the relative
    entropy and using the two partial-trace adjoint identities gives
    \[
        \re\tr
        \bigl(\omega_{XY}\log(\omega_X\otimes\omega_Y)\bigr)
        =-S(\omega_X)-S(\omega_Y).
    \]
    The asserted formula follows from
    $D(\rho\,\|\,\sigma)=-S(\rho)
        -\re\tr(\rho\log\sigma)$.
\end{proof}

\begin{theorem}[SSA equality as saturation for the product marginals]
    \label{thm:ssa_equality_product_marginal_dpi}
    \lean{isSSAEquality_iff_product_marginal_data_processing_eq}
    \lean{SSAPosDef.traceC_ABC_product_marginals}
    \leanok
    \uses{def:ssa_equality, def:quantum_relative_entropy,
        def:traceA_ABC, def:traceC_ABC, def:traceAC_ABC}
    Let $\rho_{ABC}$ be a tripartite density matrix.  Then equality holds in
    strong subadditivity if and only if relative-entropy data processing under
    the partial trace over $C$ is saturated for the pair
    $\rho_{ABC}$ and $\rho_A\otimes\rho_{BC}$:
    \[
        D(\rho_{ABC}\,\|\,\rho_A\otimes\rho_{BC})
        =D(\rho_{AB}\,\|\,\rho_A\otimes\rho_B).
    \]
    Moreover,
    \[
        \tr_C(\rho_A\otimes\rho_{BC})=\rho_A\otimes\rho_B.
    \]
    This is the product-marginal formulation in
    \cite[Equations~(5)--(7)]{Hayden2004Structure}.
\end{theorem}

\begin{proof}\leanok
    \uses{thm:relative_entropy_product_marginals}
    The two relative entropies are
    \[
        S(\rho_A)+S(\rho_{BC})-S(\rho_{ABC})
        \quad\text{and}\quad
        S(\rho_A)+S(\rho_B)-S(\rho_{AB}),
    \]
    respectively.  Cancelling the common term $S(\rho_A)$ shows that their
    equality is precisely
    \[
        S(\rho_{ABC})+S(\rho_B)=S(\rho_{AB})+S(\rho_{BC}).
    \]
    Entrywise, the partial-trace identity is
    \[
        \bigl[\tr_C(\rho_A\otimes\rho_{BC})\bigr]_{(a,b),(a',b')}
        =\sum_c(\rho_A)_{a,a'}(\rho_{BC})_{(b,c),(b',c)}
        =(\rho_A)_{a,a'}(\rho_B)_{b,b'}
        =\bigl[\rho_A\otimes\rho_B\bigr]_{(a,b),(a',b')}.
    \]
\end{proof}

\begin{theorem}[SSA equality as saturation for a maximally mixed reference]
    \label{thm:ssa_equality_maximally_mixed_reference_dpi}
    \lean{isSSAEquality_iff_maximally_mixed_reference_data_processing_eq}
    \leanok
    \uses{def:ssa_equality, def:quantum_relative_entropy,
        def:traceA_ABC, def:traceC_ABC, def:traceAC_ABC,
        lem:traceC_maximally_mixed_reference}
    Let $\rho_{ABC}$ be a tripartite density matrix and set
    \[
        \sigma_{ABC}=\frac{\mathbf 1_A}{d_A}\otimes\rho_{BC}.
    \]
    Then equality holds in strong subadditivity if and only if relative-entropy
    data processing under the partial trace over $C$ is saturated for this pair:
    \[
        D(\rho_{ABC}\,\|\,\sigma_{ABC})
        =D(\tr_C\rho_{ABC}\,\|\,\tr_C\sigma_{ABC}).
    \]
    By Lemma~\ref{lem:traceC_maximally_mixed_reference}, the reference on the
    right is $(\mathbf 1_A/d_A)\otimes\rho_B$.

    This is an equivalent hypothesis-free criterion with a different reference
    state.  The exact product-marginal formulation is
    Theorem~\ref{thm:ssa_equality_product_marginal_dpi}.
\end{theorem}

\begin{proof}\leanok
    \uses{lem:rel_entropy_eval_support, lem:kron_marginal_support,
        lem:traceC_maximally_mixed_reference}
    The two relative entropies are
    \[
        \log d_A+S(\rho_{BC})-S(\rho_{ABC})
        \quad\text{and}\quad
        \log d_A+S(\rho_B)-S(\rho_{AB}),
    \]
    respectively. The marginal support lemma places both singular references
    in the relative-entropy domain. Cancelling the common $\log d_A$ shows that
    equality of the two relative entropies is precisely
    \[
        S(\rho_{ABC})+S(\rho_B)=S(\rho_{AB})+S(\rho_{BC}).
    \]
\end{proof}

\begin{theorem}[SSA equality under factor relabeling]
    \label{thm:ssa_equality_reindex_factors}
    \lean{isSSAEquality_submatrix_prodEquiv}
    \leanok
    \uses{def:ssa_equality}
    Let $\rho$ be a Hermitian operator on $A\otimes B\otimes C$, and let
    $e_A:A'\to A$, $e_B:B'\to B$, and $e_C:C'\to C$ be bijections. Define
    $\rho'$ by
    \[
        \rho'_{(a',b',c'),(\widetilde a',\widetilde b',\widetilde c')}
        =
        \rho_{(e_A(a'),e_B(b'),e_C(c')),
          (e_A(\widetilde a'),e_B(\widetilde b'),e_C(\widetilde c'))}.
    \]
    If $\rho$ satisfies equality in strong subadditivity, then so does $\rho'$.
\end{theorem}

\begin{proof}
    \leanok
    \uses{thm:entropy_submatrix_equiv}
    The four operators entering the equality are related by the induced
    bijections:
    \[
        \rho' = \rho_{e_A\times e_B\times e_C},\qquad
        \rho'_{AB}=(\rho_{AB})_{e_A\times e_B},\qquad
        \rho'_{BC}=(\rho_{BC})_{e_B\times e_C},\qquad
        \rho'_B=(\rho_B)_{e_B}.
    \]
    These identities follow by changing variables in the three finite sums
    defining the partial traces. Entropy invariance under reindexing makes the
    corresponding four entropy terms equal, so the strong-subadditivity
    equality for $\rho$ gives the equality for $\rho'$.
\end{proof}

\begin{definition}[Quantum Markov decomposition]\label{def:hayashi_markov_decomposition}
    \lean{HayashiMarkovDecomposition}
    \leanok
    \uses{def:ssa_equality}
    A \emph{Hayashi Markov decomposition} of a tripartite state $\rho_{ABC}$
    consists of a finite direct-sum decomposition
    \[
        H_B \cong \bigoplus_j H_{B_j^L} \otimes H_{B_j^R},
    \]
    together with a unitary change of basis on $B$, a probability vector
    $(p_j)_j$, and density matrices $\rho_{A B_j^L}$ and $\rho_{B_j^R C}$ such
    that, in the adapted basis, the state becomes
    \[
        \bigoplus_j p_j\, \rho_{A B_j^L} \otimes \rho_{B_j^R C}.
    \]
    The terminology follows Hayashi's presentation of quantum Markov
    structure~\cite{Hayashi2006QuantumInformation}; the block decomposition
    used by the MPDO argument is the structure theorem of
    \cite{Hayden2004Structure}.
\end{definition}

\begin{theorem}[Hayashi SSA-equality characterization]
    \label{thm:hayashi_ssa_equality_characterization}
    \lean{hayashi_ssa_equality_characterization}
    \leanok
    \uses{def:ssa_equality, def:hayashi_markov_decomposition,
        thm:hayashi_ssa_equality_reverse}
    For a tripartite density matrix $\rho_{ABC}$,
    \[
        S(\rho_{ABC}) + S(\rho_B) = S(\rho_{AB}) + S(\rho_{BC})
    \]
    holds if and only if $\rho_{ABC}$ admits a quantum Markov decomposition on
    the middle subsystem $B$.

    The reverse implication, that a quantum Markov decomposition forces the
    equality, is proved as
    Theorem~\ref{thm:hayashi_ssa_equality_reverse}. The forward implication, that
    the equality forces the block-diagonal structure, is the deep part: it rests
    on recovery-map and Petz theory not yet available, so it is the only part
    that remains axiomatic. The biconditional combines the two.

    This equality criterion is cited here as an input for the later MPDO arguments.
    It is the equality criterion needed by Appendix~C of arXiv:1606.00608.
    The equality conditions are reviewed in
    \cite{Hayashi2006QuantumInformation, Ruskai2002Inequalities};
    the structural block-decomposition formulation used for MPDOs appears in
    \cite{Hayden2004Structure}.
\end{theorem}

\begin{theorem}[Quantum Markov structure attains SSA equality]
    \label{thm:hayashi_ssa_equality_reverse}
    \lean{hayashi_ssa_equality_characterization_reverse}
    \leanok
    \uses{def:ssa_equality, def:hayashi_markov_decomposition}
    A tripartite density matrix $\rho_{ABC}$ that admits a quantum Markov
    decomposition on the middle subsystem $B$ satisfies
    \[
        S(\rho_{ABC}) + S(\rho_B) = S(\rho_{AB}) + S(\rho_{BC}).
    \]
\end{theorem}

\begin{proof}\leanok
    \uses{def:ssa_equality, def:hayashi_markov_decomposition,
        thm:entropy_kronecker, thm:entropy_smul, thm:entropy_block_diagonal}
    Write the state in the adapted basis as the block-diagonal direct sum
    $\bigoplus_j p_j\, \rho_{A B_j^L} \otimes \rho_{B_j^R C}$. The von Neumann
    entropy of a weighted orthogonal direct sum is
    $S(\bigoplus_j p_j\, \omega_j) = -\sum_j p_j \log p_j + \sum_j p_j\,
    S(\omega_j)$ by Theorems~\ref{thm:entropy_block_diagonal}
    and~\ref{thm:entropy_smul}, and the entropy of a tensor product is additive,
    $S(\omega \otimes \tau) = S(\omega) + S(\tau)$, by
    Theorem~\ref{thm:entropy_kronecker}. Tracing out one tensor factor within each
    block, the three reduced states factor as block-diagonal direct sums,
    \[
        \rho_{AB} \cong \bigoplus_j p_j\,\rho_{AB_j^L}\otimes\rho_{B_j^R}, \qquad
        \rho_{BC} \cong \bigoplus_j p_j\,\rho_{B_j^L}\otimes\rho_{B_j^R C}, \qquad
        \rho_{B}  \cong \bigoplus_j p_j\,\rho_{B_j^L}\otimes\rho_{B_j^R},
    \]
    with $\rho_{ABC} \cong \bigoplus_j p_j\,\rho_{AB_j^L}\otimes\rho_{B_j^R C}$ itself.
    Writing $H = -\sum_j p_j\log p_j$, the four entropies expand as
    \[
        \begin{aligned}
        S(\rho_{ABC}) &= H + \sum_j p_j\bigl[S(\rho_{AB_j^L}) + S(\rho_{B_j^R C})\bigr], &
        S(\rho_{AB})  &= H + \sum_j p_j\bigl[S(\rho_{AB_j^L}) + S(\rho_{B_j^R})\bigr], \\
        S(\rho_{B})   &= H + \sum_j p_j\bigl[S(\rho_{B_j^L}) + S(\rho_{B_j^R})\bigr], &
        S(\rho_{BC})  &= H + \sum_j p_j\bigl[S(\rho_{B_j^L}) + S(\rho_{B_j^R C})\bigr].
        \end{aligned}
    \]
    Both sides of the claimed identity equal
    $2H + \sum_j p_j\bigl[S(\rho_{AB_j^L}) + S(\rho_{B_j^R C}) + S(\rho_{B_j^L}) +
    S(\rho_{B_j^R})\bigr]$, so $S(\rho_{ABC}) + S(\rho_B) = S(\rho_{AB}) +
    S(\rho_{BC})$. The basis change on $B$ and the direct-sum reindexing leave
    every entropy term unchanged.
\end{proof}
